\emph{Conventions and what is used.} The first-order theorem is Theorem \claimref{synthesis.null-plane-data-of-the-first-order-gravitational-field-are-the-shear-of-the-generators}; its clause (i) gives the current $\omega$ and its divergence, its clauses (ii) and (iv) the flux through $N$ in the gauge $h(n) = 0$ and the first-order expansion and shear, its clause (iii) the focusing equation for a field of any strength, its clause (vi) the constant $1/16\pi G$ and the first-order field equation with a source. Statements that are identities between constant-coefficient differential expressions, linear in $h$ and in $\xi$, are proved through their symbols, as in that theorem: $p$ stands for each derivative of $h$, $q$ for each derivative of $\xi$, and $h$, $\xi$ are then a constant symmetric map and a constant vector. For a bivector $\mathcal{B}$ and vectors $x, y$: $x \cdot (\mathcal{B} \cdot y) = -y \cdot (\mathcal{B} \cdot x)$, and $r \cdot (u \wedge v) = (r \cdot u)v - (r \cdot v)u$. The reciprocals of $n, \bar n$ in the frame $n, \bar n, e_{1}, e_{2}$ are $\tfrac12\bar n, \tfrac12 n$, so $\nabla = \tfrac12\bar n\,\partial_{\lambda} + \tfrac12 n\,\partial_{\bar\lambda} + \nabla_{\perp}$ with $\partial_{\lambda} = n \cdot \nabla$, $\partial_{\bar\lambda} = \bar n \cdot \nabla$, and the trace of a map is $\tfrac12\bar n \cdot h(n) + \tfrac12 n \cdot h(\bar n) + \tr_{\perp}h = n \cdot h(\bar n) + \tr_{\perp}h$.

\emph{The pair, the boost and the variations.} $B = e_{0}e_{3}$ has $B^{2} = e_{0}e_{3}e_{0}e_{3} = -e_{0}^{2}e_{3}^{2} = 1$, and $\bar n \wedge n = (e_{0} - e_{3}) \wedge (e_{0} + e_{3}) = 2\,e_{0} \wedge e_{3}$. For a vector $v$, $R_{\chi}vR_{\chi}^{\sim} = v - \chi\,B \cdot v + O(\chi^{2}) = v + \chi\,v \cdot B + O(\chi^{2})$, so the velocity is $\xi_{B}(x) = x \cdot B = (x \cdot e_{0})e_{3} - (x \cdot e_{3})e_{0}$: $e_{0} \cdot B = e_{3}$, $e_{3} \cdot B = e_{0}$, $n \cdot B = n$, $\bar n \cdot B = -\bar n$, $e_{1} \cdot B = e_{2} \cdot B = 0$. Hence $R_{\chi}nR_{\chi}^{\sim} = e^{\chi}n$, $R_{\chi}\bar nR_{\chi}^{\sim} = e^{-\chi}\bar n$, the corner is fixed pointwise, and $\xi_{B} = \lambda\,n \cdot B = \lambda n$ at $x = \lambda n + y$, $\xi_{B} = -\bar\lambda\,\bar n$ at $x = \bar\lambda\,\bar n + y$. The displacement of clause (e) of the Definition with $f(x) = x - \epsilon\,\xi(x)$ has $\underline{f}^{-1}(b) = b + \epsilon\,b \cdot \nabla\xi$ to first order, so it carries the flat field $\hb = \id$ to $\hb(a) = a + \epsilon\,a \cdot \nabla\xi$, whose derived metric is $a \cdot b - \epsilon\,\bigl((a \cdot \nabla\xi) \cdot b + a \cdot (b \cdot \nabla\xi)\bigr) = a \cdot b - \epsilon\,\Delta_{\xi}(a) \cdot b$; a rotation gauge does not change the derived metric. For the boost, the displacement by $f(x) = R_{\chi}^{\sim}xR_{\chi}$ gives $\hb(a; x) \mapsto R_{\chi}\hb(a; R_{\chi}^{\sim}xR_{\chi})R_{\chi}^{\sim}$ and the rotation gauge of rotor $R_{\chi}$ then gives $R_{\chi}\hb(R_{\chi}^{\sim}aR_{\chi}; R_{\chi}^{\sim}xR_{\chi})R_{\chi}^{\sim}$, which is the identity on the flat field and acts in the same way on the symmetric map $h$. Its derivative at $\chi = 0$: $R^{\sim}aR = a + \chi\,B \cdot a$, $R^{\sim}xR = x - \chi\,\xi_{B}$, and the outer rotor adds $-\chi\,B \cdot h(a)$; together $-\chi\bigl((\xi_{B} \cdot \nabla)h(a) + B \cdot h(a) - h(B \cdot a)\bigr) = -\chi\,L_{B}h(a)$. Also $\Delta_{\xi_{B}}(a) = a \cdot B + \nabla\bigl((x \cdot B) \cdot a\bigr) = a \cdot B - a \cdot B = 0$, since $(x \cdot B) \cdot a = -x \cdot (a \cdot B)$. The rate of change of a field along a flow, in the sense in which the null translation $\lambda \mapsto \lambda + t$ has the rate $\partial_{n}h$ in the first-order theorem, is its derivative along the flow: $\Delta_{\xi}$ for the flat field and $L_{B}h$ for the exact boost. $L_{B}$ is that derivative for the first-order field, so it commutes with $\mathcal{E}$ and carries displacements to displacements, $L_{B}\Delta_{\zeta} = \Delta_{[\xi_{B}, \zeta]}$ (verified in the check on polynomial fields).

\emph{(i).} In symbols $\Delta_{\xi}(a) = (a \cdot q)\xi + (a \cdot \xi)q =: k(a)$, with trace $2\,q \cdot \xi$, divergence $k(q) = Q(q)\xi + (q \cdot \xi)q$, and $\ip{h}{k} = 2\,q \cdot h(\xi)$. Insert these in the current of clause (i) of the first-order theorem, with $d = h(p)$:
\[
\omega(h, k) = \tfrac12(p - q)\ip{h}{k} + \tfrac12\bigl(\tau\,k(p) - 2(q \cdot \xi)h(q)\bigr) + \tfrac12\bigl(2(q \cdot \xi)h(p) - \tau\,k(q)\bigr) + h(k(q)) - k(h(p)) - \tau\,(q \cdot \xi)(p - q) .
\]
With $k(p) = (p \cdot q)\xi + (p \cdot \xi)q$, $h(k(q)) = Q(q)h(\xi) + (q \cdot \xi)h(q)$ and $k(h(p)) = (q \cdot h(p))\xi + (\xi \cdot h(p))q$ the two terms in $h(q)$ cancel and
\[
\omega(h, k) = (p - q)\,(q \cdot h(\xi)) + (q \cdot \xi)h(p) + Q(q)h(\xi) - (q \cdot h(p))\xi - (\xi \cdot h(p))q + \tau\bigl(\tfrac12(p \cdot q)\xi + \tfrac12(p \cdot \xi)q - \tfrac12 Q(q)\xi + \tfrac12(q \cdot \xi)q - (q \cdot \xi)p\bigr) .
\]
The symbol of $\mathcal{E}$ is the stored Einstein map, written out in the proof of the first-order theorem: $2\,\mathcal{E}[h](\xi) = Q(p)h(\xi) + (\xi \cdot p)\tau\,p - (\xi \cdot p)h(p) - (\xi \cdot h(p))p - \xi\,(Q(p)\tau - p \cdot h(p))$. The symbol of $\Psi_{\xi}[h]$ is $\xi \wedge (h(p) - \tau p) + (q - p) \wedge h(\xi) - \tfrac12\tau\,q \wedge \xi$, and that of its divergence is its inner product with $r = p + q$:
\[
r \cdot \Psi = (r \cdot \xi)h(p) - (r \cdot h(p))\xi - \tau(r \cdot \xi)p + \tau(r \cdot p)\xi + (Q(q) - Q(p))h(\xi) - (r \cdot h(\xi))(q - p) - \tfrac12\tau(r \cdot q)\xi + \tfrac12\tau(r \cdot \xi)q .
\]
Add $2\,\mathcal{E}[h](\xi)$ and collect. The terms in $h(\xi)$: $Q(p) + Q(q) - Q(p) = Q(q)$. In $h(p)$: $-(\xi \cdot p) + (p \cdot \xi) + (q \cdot \xi) = q \cdot \xi$. In $\xi$ without $\tau$: $p \cdot h(p) - p \cdot h(p) - q \cdot h(p) = -q \cdot h(p)$. In $p$ and $q$ without $\tau$, using $p \cdot h(\xi) = \xi \cdot h(p)$: $-(\xi \cdot h(p))p + (p \cdot h(\xi))(p - q) + (q \cdot h(\xi))(p - q) = -(\xi \cdot h(p))q + (p - q)(q \cdot h(\xi))$. In $\tau$: $(\xi \cdot p)p - Q(p)\xi - (p \cdot \xi)p - (q \cdot \xi)p + Q(p)\xi + (p \cdot q)\xi - \tfrac12(p \cdot q)\xi - \tfrac12 Q(q)\xi + \tfrac12(p \cdot \xi)q + \tfrac12(q \cdot \xi)q$, which is the bracket above. So $2\,\mathcal{E}[h](\xi) + r \cdot \Psi = \omega(h, k)$ for all $h, \xi, p, q$, which is the identity. The check verifies it completely: both sides are linear in $h$ and in $\xi$ and homogeneous of degree two in the pair $(p, q)$, and a homogeneous quadratic polynomial in eight variables that vanishes at the eight unit points and at the twenty-eight sums of two of them vanishes identically; the check evaluates the difference at those $36$ points for the ten unit symmetric maps and the four frame vectors, $1440$ evaluations, using the stored curvature formula for $\mathcal{E}$ and not the formula above. Uniqueness: the symbols of the six listed bivectors are $\xi \wedge h(p)$, $\tau\,\xi \wedge p$, $p \wedge h(\xi)$, $q \wedge h(\xi)$, $\tau\,q \wedge \xi$, $h(q) \wedge \xi$; on four generic symbols the sixteen components of their divergences have rank six (computed), so the coefficients are determined, and solving that system against $\omega(h, \Delta_{\xi}) - 2\,\mathcal{E}[h](\xi)$ returns $1, -1, -1, 1, -\tfrac12, 0$ (computed). \emph{The six are all.} A bivector $\mathcal{K}(h, \xi, p)$ linear in a symmetric map and in two vectors commutes with the rotors, $\mathcal{K}(Rh(\tilde R\,\cdot\,R)\tilde R, R\xi\tilde R, Rp\tilde R) = R\,\mathcal{K}(h, \xi, p)\,\tilde R$, exactly when the form $(u \wedge v) \cdot \mathcal{K}(h, \xi, p)$ is annihilated by each unit bivector acting on every vector by the commutator product, as in the proof of clause (v) of the first-order theorem. Its coefficients on (blade, symmetric pair, frame vector, frame vector) number $6 \cdot 10 \cdot 4 \cdot 4 = 960$, and the six unit bivectors give a linear system with integer coefficients whose nullity modulo the prime $1000003$ is six (computed); a minor that is nonzero modulo a prime is nonzero, so the rational space of solutions has dimension at most six. The three bivectors $\xi \wedge h(p)$, $\tau\,\xi \wedge p$, $p \wedge h(\xi)$ are built with the product, so they commute with the rotors, and so do their duals, $I_{4}$ times them, since $I_{4}$ commutes with the rotors; the six are independent (computed on four generic arguments). So the space has dimension six. Under the reflection $v \mapsto e_{1}ve_{1}$ of the frame the three are carried along, $\mathcal{K}(h', \xi', p') = \mathcal{K}(h, \xi, p)'$, and their duals change sign, because $I_{4}$ anticommutes with $e_{1}$ (computed on the same arguments); so the bivectors of this kind that are even under reflection are the combinations of the three. With the derivative placed on $h$ ($p$) or on $\xi$ ($q$) that gives the six listed. The last sentence of the clause: if $\mathcal{E}[h] = 0$ then $\omega(h, \Delta_{\xi}) = \nabla \cdot \Psi_{\xi}[h]$, and the divergence of the divergence of a bivector vanishes.

\emph{(ii).} $\Psi_{\xi}[h]$ contains $\xi$ and its first derivatives only. Where $\xi = 0$, $\Psi_{\xi}[h] = \dot\nabla \wedge h(\dot\xi) - \tfrac12\tau\,\nabla \wedge \xi = e^{\mu} \wedge h(e_{\mu} \cdot \nabla\xi) - \tfrac12\tau\,e^{\mu} \wedge (e_{\mu} \cdot \nabla\xi)$, summed over $\mu$. With $e_{\mu} \cdot \nabla\xi = w\,e_{\mu} \cdot B$ on $C$: only $\mu = 0, 3$ contribute, $e^{0} = e_{0}$, $e^{3} = -e_{3}$, so $\dot\nabla \wedge h(\dot\xi) = w\,(e_{0} \wedge h(e_{3}) - e_{3} \wedge h(e_{0}))$ and $\nabla \wedge \xi = w\,(e_{0} \wedge e_{3} - e_{3} \wedge e_{0}) = 2wB$. The scalar part of a product of two blades is $(a \wedge b) \cdot (c \wedge d) = (a \cdot d)(b \cdot c) - (a \cdot c)(b \cdot d)$, so $B \cdot B = 1$, $B \cdot (e_{0} \wedge h(e_{3})) = -e_{3} \cdot h(e_{3})$ and $B \cdot (-e_{3} \wedge h(e_{0})) = e_{0} \cdot h(e_{0})$. Hence $B \cdot \Psi_{\xi}[h] = w\,(e_{0} \cdot h(e_{0}) - e_{3} \cdot h(e_{3}) - \tau)$; the first two terms are the trace over the plane of $n, \bar n$, $e^{0} \cdot h(e_{0}) + e^{3} \cdot h(e_{3}) = n \cdot h(\bar n)$, so $B \cdot \Psi_{\xi}[h] = -w\,\tr_{\perp}h$. No property of $h$ was used. On $C$, $\Delta_{\xi}(a) = w\,(a \cdot B + \nabla((x \cdot B) \cdot a)) + (a \cdot \nabla w)\xi_{B} + \nabla w\,(\xi_{B} \cdot a) = 0$ for $\xi = w\,\xi_{B}$, and in general $\Delta_{\xi} = w\,\Delta_{\xi_{B}} = 0$ there because it is linear in the first derivatives of $\xi$, which are those of $w\,\xi_{B}$ on $C$; so a term $\nabla \cdot Y(h, \Delta_{\xi})$ added to the pairing changes the flux through a cut by $Y(h, \Delta_{\xi})$ on it, which vanishes on $C$. For a displacement, $\tr_{\perp}\Delta_{\zeta} = \sum_{j = 1, 2}e^{j} \cdot \bigl(e_{j} \cdot \nabla\zeta + \nabla(\zeta \cdot e_{j})\bigr) = 2\sum_{j}e^{j} \cdot \partial_{j}\zeta = 2\,\nabla_{\perp} \cdot \zeta_{\perp}$, which involves only derivatives along $C$ of the part of $\zeta$ along $C$; it vanishes on $C$ when that part vanishes on $C$. The check verifies the corner value on a generic polynomial field in no gauge with a generic weight $w$, that it is not the full trace, the formula for $\tr_{\perp}\Delta_{\zeta}$, the two kinds of displacement that leave the value unchanged, and for a displacement along $C$ that the total over a window equals the flux of $-2\zeta_{\perp}$ through its edge.

\emph{(iii).} By clause (a) of the Definition the physical vector over a background vector $a$ is $\hb^{-1}(a)$, compared with the flat product; the physical bivector over $e_{1} \wedge e_{2}$ is the outer product of the physical vectors, and its size is the area that physical rods measure on the element of the corner of unit background area. For vectors $u, v$, $(u \wedge v)^{2} = (u \cdot v)^{2} - Q(u)Q(v)$, which gives the second form with the derived metric of clause (d). A rotation gauge turns both physical vectors by one rotor (clause (e)) and leaves the size unchanged. To first order $g(e_{j}, e_{l}) = e_{j} \cdot e_{l} + h(e_{j}) \cdot e_{l}$, with $e_{j} \cdot e_{l} = -\delta_{jl}$ on the corner, so $\mathsf{a}^{2} = (1 - h(e_{1}) \cdot e_{1})(1 - h(e_{2}) \cdot e_{2}) + O(h^{2}) = 1 + \tr_{\perp}h + O(h^{2})$, as $e^{j} = -e_{j}$. The check expands $\mathsf{a}^{2}$ exactly as a polynomial in $\epsilon$ for $\hb^{-1} = \id + \epsilon s$ on four generic symmetric maps in no gauge, and confirms $-(g_{1} \wedge g_{2})^{2} = \mathsf{a}^{2}$ on a field that is not weak.

\emph{(iv).} For a bivector field $\Psi$ and the constant vector $n$, $n \cdot (\nabla \cdot \Psi) = -\nabla \cdot (n \cdot \Psi) = \nabla \cdot (\Psi \cdot n)$. The vector $J = \Psi \cdot n$ has $n \cdot J = 0$, so $J = \alpha\,n + J_{\perp}$ with $\alpha = \tfrac12\bar n \cdot J$; writing $\Psi = \kappa B + \dots$ with $\kappa = B \cdot \Psi$, the other blades of the frame $n, \bar n, e_{1}, e_{2}$ contribute nothing to $\bar n \cdot (\Psi \cdot n)$ and $(\kappa B) \cdot n = -\kappa\,n \cdot B = -\kappa n$, so $\alpha = -B \cdot \Psi$. With the form of $\nabla$ above, $\nabla \cdot (\alpha n) = \tfrac12(\bar n \cdot n)\partial_{\lambda}\alpha = \partial_{\lambda}\alpha$. Likewise $\bar n \cdot (\nabla \cdot \Psi) = \nabla \cdot (\Psi \cdot \bar n)$, $(\kappa B) \cdot \bar n = -\kappa\,\bar n \cdot B = \kappa\bar n$, and the part along $\bar n$ is $\partial_{\bar\lambda}(B \cdot \Psi)$. Integrate $n \cdot \omega(h, \Delta_{\xi}) = 2\,n \cdot \mathcal{E}[h](\xi) + n \cdot (\nabla \cdot \Psi_{\xi}[h])$ over $W$: the second term gives $\int_{U}(-B \cdot \Psi)\big|_{\Lambda} - \int_{U}(-B \cdot \Psi)\big|_{0}$ and the flux of $J_{\perp}$ through the sides, and by (ii) and (iii) $-B \cdot \Psi = w\,\tr_{\perp}h = 2w\,\mathsf{a}_{1}$ on $C$. If $\xi$ and its first derivatives vanish on the far cut and on the sides, $\Psi_{\xi}[h]$ vanishes there. On $\bar N_{+}$ the corner is the lower limit of $\partial_{\bar\lambda}(+B \cdot \Psi)$, which gives $-(B \cdot \Psi)\big|_{C} = +2w\,\mathsf{a}_{1}$. The measure on $\bar N$ is $d\bar\lambda\,d^{2}y$ as on $N$, by the exchange of $n$ and $\bar n$. For the pair, the corner terms of the two halves are $-\tfrac{1}{8\pi G}\int w_{N}\mathsf{a}_{1}$ and $+\tfrac{1}{8\pi G}\int w_{\bar N}\mathsf{a}_{1}$. The generator: by definition $\delta H_{X} = \tfrac{1}{16\pi G}\Phi(\delta h, X)$ on each half, and the generator of the pair is defined as the sum over the two halves (the light-front setting: the flux through $N_{+}$ together with $\bar N_{+}$ taken as the symplectic form of the pair; not proved). The rate of the flow is $X = \Delta_{\xi}$ plus terms of first order in $h$; the first part does not depend on $h$, so the part of the generator linear in $h$ is the displayed corner term, the first-order field $h$ standing for the variation, and the terms of first order in $h$ give a part quadratic in $h$, which for the exact boost is the generator of (vi)(a). For the exact boost of both halves $\Delta_{\xi_{B}} = 0$ and the linear part is zero for every $h$, a solution or not. The relative boost at unit rate, $w_{N} = 0$, $w_{\bar N} = 1$, is the flow $x \mapsto R_{\chi}xR_{\chi}^{\sim}$ near $C$ on $\bar N_{+}$ and the identity on $N_{+}$: it takes $\bar n$ to $e^{-\chi}\bar n$ and leaves $n$. The orientation is that of the first-order theorem, in which the generator $P$ of $\lambda \mapsto \lambda + t$ is positive: $\tfrac{1}{16\pi G}\Phi(h', \partial_{n}h)$ is, on the two components and after an integration by parts along the generators, $\tfrac{1}{16\pi G}\int\ip{\partial_{n}h'^{TT}}{\partial_{n}h^{TT}} = dP(h')$. The check takes a polynomial solution of $\mathcal{E}[h] = 0$ in a generic gauge, the sum of powers of $p \cdot x$ at three null momenta that cross $N$, each times a map in the kernel of the stored Einstein map at that momentum (the six-dimensional kernel of Theorem \claimref{synthesis.first-order-gravitational-field-has-two-free-directions-at-null-momentum-and-four-constraints}: a transverse trace-free part and a displacement), plus a polynomial displacement, and a field $\xi = f g\,\xi_{B}$ with $f$ vanishing to second order on the far cut and $g$ on the sides; it integrates $n \cdot \omega(h, \Delta_{\xi})$ exactly over the window and finds $-2\int g\,\mathsf{a}_{1}$ on the corner, with no remainder, and $+2\int g\,\mathsf{a}_{1}$ through the window of $\bar N_{+}$ with the flow cut off on that half. A second check does the same for a solution that is not built from plane waves: a generic homogeneous quadratic field, for which $\mathcal{E}$ is a constant map, less a combination of fields with the same image under $\mathcal{E}$ found by exact elimination; it is not a displacement (a component of its first-order curvature is nonzero) and not in the gauge; and it finds off the solutions the remainder $2\int n \cdot \mathcal{E}[h](\xi)$.

\emph{(v).} With $\Delta_{\xi_{B}} = 0$ the identity of (i) is $\nabla \cdot \Psi_{\xi_{B}}[h] = -2\,\mathcal{E}[h](\xi_{B})$; contracted with $n$ on $N$, where $\xi_{B} = \lambda n$, and split as in (iv), it is the displayed law. In the gauge $h(n) = 0$ on $N$: $h(\xi_{B}) = \lambda h(n) = 0$ and $(\partial_{\mu}h)(n) = 0$, so the term $\dot\nabla \wedge \dot h(\xi_{B})$ vanishes; $h(e_{3}) = -h(e_{0})$, so $\dot\nabla \wedge h(\dot\xi_{B}) = e_{0} \wedge h(e_{3}) - e_{3} \wedge h(e_{0}) = -n \wedge h(e_{0})$ and $B \cdot (-n \wedge h(e_{0})) = n \cdot h(e_{0}) = e_{0} \cdot h(n) = 0$; $B \cdot (\lambda\,n \wedge (d - \nabla\tau)) = -\lambda\,n \cdot (d - \nabla\tau) = \lambda\,\partial_{n}\tau$, since $n \cdot d = \dot h(n) \cdot \dot\nabla = 0$; and $-\tfrac12\tau\,B \cdot (\nabla \wedge \xi_{B}) = -\tau$. So $-B \cdot \Psi_{\xi_{B}}[h] = \tau - \lambda\,\partial_{n}\tau$ and $\partial_{\lambda}$ of it is $-\lambda\,\partial_{n}^{2}\tau$. By clause (ii) of the first-order theorem $n \cdot \mathcal{E}[h](n) = \tfrac12\partial_{n}^{2}\tau$ in the gauge, so the right side of the law is $-\lambda\,\partial_{n}^{2}\tau$ too, and the transverse term vanishes. By its clause (iv), $\theta = \tfrac12\partial_{n}\tau$ and the equation along the generators is the first-order focusing equation $\partial_{n}\theta = n \cdot \mathcal{E}[h](n)$; multiplied by $\lambda$ and integrated, $\int_{0}^{\Lambda}\lambda\,\partial_{\lambda}\theta = \Lambda\theta(\Lambda) - \tfrac12(\tau(\Lambda) - \tau(0))$, which is the same statement. In the gauge the trace over the corner plane is the whole trace ($n \cdot h(\bar n) = \bar n \cdot h(n) = 0$), so $\mathsf{a}_{1} = \tfrac12\tau$ along the whole generator, each cut $\lambda = \text{const}$ of $N$ being a plane parallel to $C$ with the same area element by (iii). With $n \cdot \mathcal{E}[h](n) = -8\pi G\,n \cdot T(n)$: $-8\pi G\int_{0}^{\Lambda}\lambda\,n \cdot T(n) = \Lambda\,\partial_{\lambda}\mathsf{a}_{1}(\Lambda) - \mathsf{a}_{1}(\Lambda) + \mathsf{a}_{1}(0)$. In any gauge the law is integrated over $W$ as in (iv). About a cut $\lambda = \lambda_{0}$: the same integration by parts with the weight $\lambda - \lambda_{0}$ from $\lambda_{0}$ gives $\int_{\lambda_{0}}^{\Lambda}(\lambda - \lambda_{0})\,\partial_{\lambda}\theta = (\Lambda - \lambda_{0})\theta(\Lambda) - \tfrac12(\tau(\Lambda) - \tau(\lambda_{0}))$; equivalently the field translated along $n$ by $\lambda_{0}$ is again a first-order field of the gauge with the translated source. The check verifies the conservation law and its two splittings as polynomial identities on a generic field with a source in no gauge, the balance on a window with the side fluxes by exact integration, and in the gauge the three statements just proved and the equality of the two routes as numbers.

\emph{(vi)(a).} For transverse $a, b$: $B \cdot a = 0$ and $(B \cdot h(a)) \cdot b = -h(a) \cdot (B \cdot b) = 0$, so $L_{B}h(a) \cdot b = (\xi_{B} \cdot \nabla)h(a) \cdot b$; and $L_{B}h(n) = (\xi_{B} \cdot \nabla)h(n) + B \cdot h(n) - h(B \cdot n) = h(n) = 0$ in the gauge, since $B \cdot n = -n$. The trace of $a \mapsto B \cdot h(a) - h(B \cdot a)$ vanishes, so $\tr L_{B}h = \xi_{B} \cdot \nabla\tau$. On $N$, where $\xi_{B} = \lambda n$ and $\partial_{n}\xi_{B} = n \cdot B = n$: $(L_{B}h)^{TT} = \lambda\,\partial_{n}h^{TT}$, $\partial_{n}(L_{B}h)^{TT} = \partial_{n}h^{TT} + \lambda\,\partial_{n}^{2}h^{TT}$, and the same for the trace. Insert in the flux formula of clause (ii) of the first-order theorem, which holds for two fields of the gauge:
\[
n \cdot \omega(h, L_{B}h) = \tfrac12\bigl(\lambda\ip{\partial_{n}h^{TT}}{\partial_{n}h^{TT}} - \ip{h^{TT}}{\partial_{n}h^{TT}} - \lambda\ip{h^{TT}}{\partial_{n}^{2}h^{TT}}\bigr) - \tfrac14\bigl(\lambda(\partial_{n}\tau)^{2} - \tau\,\partial_{n}\tau - \lambda\,\tau\,\partial_{n}^{2}\tau\bigr) .
\]
Since $\partial_{\lambda}(\lambda\ip{h^{TT}}{\partial_{n}h^{TT}}) = \ip{h^{TT}}{\partial_{n}h^{TT}} + \lambda\ip{\partial_{n}h^{TT}}{\partial_{n}h^{TT}} + \lambda\ip{h^{TT}}{\partial_{n}^{2}h^{TT}}$ and likewise for $\tau$, this is $\lambda\ip{\partial_{n}h^{TT}}{\partial_{n}h^{TT}} - \tfrac12\lambda(\partial_{n}\tau)^{2} - \tfrac12\partial_{\lambda}\bigl(\lambda\ip{h^{TT}}{\partial_{n}h^{TT}} - \tfrac12\lambda\,\tau\,\partial_{n}\tau\bigr)$. With $\sigma = \tfrac12\partial_{n}h^{TT}$, $\theta = \tfrac12\partial_{n}\tau$ the first two terms are $4\lambda(\ip{\sigma}{\sigma} - \tfrac12\theta^{2})$, and $\ip{h}{\partial_{n}h} = \ip{h^{TT}}{\partial_{n}h^{TT}} + \tfrac12\tau\,\partial_{n}\tau$ for fields of the gauge (the decomposition in the proof of that clause), which gives the displayed form. The term under $\partial_{\lambda}$ carries the factor $\lambda$ and vanishes on $C$. The form $\ip{\sigma}{\sigma} - \tfrac12\theta^{2}$ is not positive definite: $\ip{\sigma}{\sigma} \geq 0$ is the part of the two components and $-\tfrac12\theta^{2}$ that of the trace; on solutions in the gauge $\partial_{n}^{2}\tau = 0$ (clause (ii) of the first-order theorem), so $\theta$ is constant along each generator, and it need not vanish. The check exhibits a field of the gauge for which the integral over the window is negative. The generator: the current is covariant, so $n \cdot \omega(L_{B}h, k) + n \cdot \omega(h, L_{B}k) = \partial_{\lambda}(\lambda\,n \cdot \omega(h, k))$ on $N$ (the derivative of the vector $\omega$ along $\xi_{B}$, contracted with $n$, is $\xi_{B} \cdot \nabla(n \cdot \omega) + n \cdot \omega$; verified in the check); hence $\delta\bigl(\tfrac{1}{32\pi G}\Phi_{W}(h, L_{B}h)\bigr) = \tfrac{1}{16\pi G}\Phi_{W}(\delta h, L_{B}h)$ up to a term on the far cut, which is the definition of the generator of the flow with rate $L_{B}h$. That the flux through $N$ is the whole symplectic form is the light-front setting of the first-order theorem, taken as there.

\emph{(vi)(b).} By clause (iv) of the first-order theorem, in this gauge $\hb^{-1}(n) = n$, the generators of $N$ run along $K = n$ with $K \cdot D = \partial_{n}$ on scalars, and $\tfrac12\partial_{n}g(a, b) = \hb^{-1}(a) \cdot \mathsf{B}(\hb^{-1}(b))$ for background $a, b$, with $\mathsf{B}(A) = A \cdot D\,K$ of its clause (iii). Put $g_{j} = \hb^{-1}(e_{j})$, $j = 1, 2$. Then $g_{j} \cdot n = (e_{j} + s(e_{j})) \cdot n = e_{j} \cdot s(n) = 0$, and $g_{1}, g_{2}$ are independent, so they represent the screen of $n$. Let $G$ be the Gram matrix of $g_{1}, g_{2}$, with entries $g_{j} \cdot g_{l}$, and $S$ the matrix of the screen map $\mathsf{S}$ in the basis $g_{1}, g_{2}$ ($\mathsf{B}(g_{l})$ is the corresponding combination of $g_{1}, g_{2}$ modulo $n$, and $g_{j} \cdot n = 0$). The relation reads $\tfrac12\partial_{n}G = GS$ as matrices, so $\theta = \tr\mathsf{S} = \tfrac12\tr(G^{-1}\partial_{n}G) = \tfrac12\partial_{n}\ln\det G$, and $\det G = \mathsf{a}^{2}$ by (iii). So $\theta = \partial_{\lambda}\ln\mathsf{a}$, $\partial_{\lambda}^{2}\mathsf{a} = \mathsf{a}(\partial_{\lambda}\theta + \theta^{2})$, and the focusing equation of clause (iii)(c) of that theorem, $\partial_{\lambda}\theta = -\tfrac12\theta^{2} - \ip{\sigma}{\sigma} + n \cdot \mathcal{R}(n)$, gives the statement. The check verifies $\theta = \tfrac12\partial_{n}(\mathsf{a}^{2})/\mathsf{a}^{2}$ and the second derivative of $\mathsf{a}^{2}$ exactly on the field of that theorem's check, which is not weak.

\emph{(vi)(c).} $\partial_{\lambda}(\lambda\,\partial_{\lambda}\mathsf{a} - \mathsf{a}) = \lambda\,\partial_{\lambda}^{2}\mathsf{a}$; integrate from $0$ to $\Lambda$ and insert (b) and (F). For the family, $\mathsf{a} = 1 + \epsilon\mathsf{a}_{1} + \epsilon^{2}\mathsf{a}_{2} + \dots$ with $\mathsf{a}_{1} = \tr_{\perp}s_{1} = \tfrac12\tr_{\perp}h$ for $h = 2s_{1}$, by (iii); $\theta$ and $\sigma$ are of order $\epsilon$, with coefficients of $\epsilon$ equal to $\tfrac12\partial_{n}\tau$, $\tfrac12\partial_{n}h^{TT}$ (clause (iv) of the first-order theorem), and $\ip{\sigma}{\sigma}$ on the screen is at order $\epsilon^{2}$ the pairing of maps. The integrand $\lambda\,\mathsf{a}\,(\ip{\sigma}{\sigma} - \tfrac12\theta^{2} + 8\pi G\,n \cdot T(n))$ is $8\pi G\,\lambda\,n \cdot T_{1}(n)$ at order $\epsilon$, which is (v), and $\lambda\,(\ip{\sigma}{\sigma} - \tfrac12\theta^{2} + 8\pi G\,n \cdot T_{2}(n) + 8\pi G\,\mathsf{a}_{1}\,n \cdot T_{1}(n))$ at order $\epsilon^{2}$. Without matter the right side at order $\epsilon^{2}$ is $8\pi G$ times $\tfrac{1}{8\pi G}\int\lambda(\ip{\sigma}{\sigma} - \tfrac12\theta^{2})$, the quantity of (a). Hypothesis (F) is the field equation of Theorem \claimref{synthesis.field-equation-linear-in-the-curvature-with-a-conserved-source-is-einsteins}(vii) contracted with $n$ twice: $n \cdot \mathcal{G}(n) = n \cdot \mathcal{R}(n)$ and $n \cdot (\Lambda n) = 0$ for $Q(n) = 0$; it is stated as a hypothesis and that theorem's conditions are not assumed here. This clause has no separate check: it follows from (a) and (b) by the integration and the expansion just given.

\emph{(vii)(a).} $I_{4}B = e_{0}e_{1}e_{2}e_{3}e_{0}e_{3} = -e_{0}e_{1}e_{2}e_{0}e_{3}e_{3} = e_{0}e_{1}e_{2}e_{0} = e_{1}e_{2}$, since $e_{3}^{2} = -1$ and $e_{0}$ commutes with $e_{1}e_{2}$. $B^{*} = e_{1}e_{2}$ squares to $-1$ and commutes with $B$ and with $n, \bar n$. The continuation is clause (i) of Theorem \claimref{synthesis.modular-period-on-values-is-the-dual-plane-turn}, with its $B$ and $S(\lambda)$ the present $B$ and $R_{\chi}$; its clause (ii) records that on points the continuation is not this rotor. Conjugation by $e^{-\varphi B^{*}/2}$ fixes $n$ and $\bar n$ and turns $e_{1}, e_{2}$ by $\varphi$; a map $a \mapsto (a \cdot n)m + (a \cdot m)n$ with $m$ in the corner is carried to the same map with $m$ turned, the maps $\epsilon_{+}, \epsilon_{\times}$ turn by $2\varphi$ (clause (viii) of the first-order theorem), and the maps built on $n, \bar n$ alone and the projection on the corner are unchanged. $e^{-\pi B^{*}/2} = -B^{*}$ and $e^{-2\pi B^{*}/2} = -1$. The check verifies these statements with exact rotors.

\emph{(vii)(b).} Clause (i) of the imported Fact \claimref{synthesis.imported-null-plane-modular-hamiltonian-first-law} gives the modular Hamiltonian $2\pi K(0)$ of the vacuum on the algebra of the half $u > 0$ of a null plane, $K(0) = \int d^{2}y\int_{0}^{\infty}u\,T_{uu}\,du$, with $u$ affine along the generators; here $u = \lambda$, $\partial_{u} = n \cdot \nabla$ and $T_{uu} = n \cdot T(n)$. Its clause (iii) defines $\Delta S = \Delta\langle 2\pi K\rangle - S(\rho\,\|\,\omega_{0})$, with $\Delta\langle K\rangle = \int\!\!\int\lambda\,n \cdot (\rho(T) - \omega_{0}(T))(n)$. Under (S) the first-order field obeys the equation of (v) with this source, which is supported in the window, so by (v), integrated over $U$, $\Delta\langle K\rangle = \tfrac{1}{8\pi G}\int_{U}\bigl((\mathsf{a}_{1} - \lambda\,\partial_{\lambda}\mathsf{a}_{1})|_{\Lambda} - \mathsf{a}_{1}|_{0}\bigr)$, and $2\pi/8\pi G = 1/4G$. The relative entropy is nonnegative and of second order in $\rho - \omega_{0}$ (clause (iii) of the Fact); (S) asks that it and $\Delta\langle K\rangle$ be finite, as that clause does. About another cut the Fact's clause (i) gives $2\pi K(u_{0})$ and (v) the area combination about the cut. Nothing here is derived about the state: the period $2\pi$ is the imported Fact's. The check forms $2\pi K(u_{0})$ for a polynomial source about two cuts and compares it with the area combination over $4G$; that computes the first law about a cut and the factor $2\pi/8\pi G$, and nothing thermal.
